\paragraph{A constant-penalty limit compatible with multiple zeros.}
The earlier pointwise spatial comparison fails at a
double-zero collision. We replace it by a division-free
\emph{heat-convolved positive defect} without assuming
a sign for the unknown Laguerre expression.
Let $T=1/2$, $k_n=(n+1)(2n+1)$, and
$P_h=\exp(\frac h2\partial_x^2)$.
For $C>0$, define
\begin{align}
E_{n,C}(r,x)
&=[k_nU_{n+1}(r,x)-CU_n(r,x)]_+,
\label{eq:penalty-defect}\\
\mathcal E_{n,C}(s,x)
&=\int_0^s(P_{s-r}E_{n,C}(r,\cdot))(x)\,dr.
\label{eq:penalty-convolution}
\end{align}
Unlike a ratio, these expressions remain defined
even at a multiple real zero.

\begin{theorem}[Constant-penalty heat criterion]
\label{thm:penaltycriterion}
For the genuine theta heat-flow family
$U_n(s,x)=L_{n,1/2-s}(x)$ the following are
equivalent:
\begin{enumerate}
\item the Riemann hypothesis;
\item $\lim_{C\to\infty}\mathcal E_{n,C}(T,x)=0$
for every $n\ge1$ and real $x$;
\item $U_n(s,x)\ge0$ for all $n\ge1$,
$0\le s\le T$ and real $x$.
\end{enumerate}
If $U_n$ is negative anywhere at positive
backward time, the penalty convolution grows
at least linearly with $C$ at every spatial
point for a suitable later time.
\end{theorem}

\begin{proof}
For each fixed $s$, $x\mapsto U_n(s,x)$
is real analytic and not identically zero.
To see the latter directly from the genuine
theta integral, use
$H_t(iy)=\int_0^\infty e^{tu^2}\Phi(u)
\cosh(yu)\,du$. Its even Taylor coefficients
are all positive. Those of
$H_t(iy)^2=\sum_nL_{n,t}(0)y^{2n}$
are likewise strictly positive. Thus
$U_n(s,0)>0$, and its real zero set
has measure zero for every fixed $s$.

By the de Bruijn--Newman theorem and
Rodgers--Tao's lower bound, RH implies
$\Lambda=0$; hence all transforms
$H_t$, $t\ge0$, have only real zeros.
Their generalized Laguerre
inequalities prove (1)$\Rightarrow$(3).
The assertion (3) at $s=T$ conversely
gives RH by the classical all-order
Laguerre--P\'olya characterization.

Under (3), one has
$0\le E_{n,C}\le k_n[U_{n+1}]_+$,
and the right side is bounded on the
compact time slab. At every point where
$U_n>0$, the defect is eventually zero.
For each time the remaining zeros of
$U_n$ form a measure-zero set in space.
Fubini's theorem and dominated
convergence in the Gaussian heat kernel
prove (3)$\Rightarrow$(2).

The heat hierarchy and the definition
of $E_{n,C}$ give the differential
inequality
\[
(\partial_s-\tfrac12\partial_x^2+C)U_n
\ge-E_{n,C}.
\]
Since $U_n(0,\cdot)\ge0$, positivity
of the heat semigroup implies
\begin{align}
U_n(s,x)
&\ge-\int_0^s e^{-C(s-r)}
(P_{s-r}E_{n,C}(r,\cdot))(x)\,dr
\nonumber\\
&\ge-\mathcal E_{n,C}(s,x).
\label{eq:penalty-barrier}
\end{align}
Assertion (2) at $s=T$ implies
$U_n(T,x)\ge0$, proving RH.

For the final statement, a negative
value at positive backward time
implies by continuity a nonempty
space--time rectangle on which
$U_n\le-\delta<0$.
Since $U_{n+1}$ is locally bounded,
$E_{n,C}\ge C\delta-O(1)$ there.
The strictly positive Gaussian
kernel gives that rectangle
positive weight at any subsequent
time, so the convolution grows
at least linearly in $C$.
\end{proof}

\begin{proposition}[Sharp collision and negative-region tests]
\label{prop:penalty-collision}
For the backward heat polynomial
$H_t(z)=z^2-2t$, write $s=T-t$
and let $T=1/2$. Then
\[
U_1(s,x)=2x^2+2-4s,\qquad U_2=1.
\]
At the endpoint double zero,
$E_{1,C}(T,0)=6$ for every $C>0$,
but, for $C\ge3$ and every real $x$,
\begin{equation}\label{eq:penalty-collision-rate}
0\le\mathcal E_{1,C}(T,x)\le9/C.
\end{equation}
Moreover the unweighted Duhamel
initial contribution and source
are both exactly $3$ at $(T,0)$.
In contrast, the negative target
$U(s,x)=1-2s$, with $U_2=1/3$
and $k_1=6$, obeys the same heat
hierarchy and nonnegative initial
data, but has
\begin{equation}
\mathcal E_C(1,x)\ge3C/16
\quad(C>0,\ x\in\mathbb R).
\label{eq:penalty-negative}
\end{equation}
\end{proposition}
\begin{proof}
Since $U_1(r,y)\ge2-4r$ on
$0\le r\le T$, the defect can
be nonzero only when
$r>T-3/(2C)$. It is everywhere
at most $6$. Gaussian convolution
preserves the upper bound, and
integration over this time window
gives $9/C$. At $(T,0)$ the
defect equals $6$. Since
$P_T(x^2)=x^2+T$,
$P_TU_1(0,\cdot)(0)=3=6T$.

In the negative model,
$E_C(r,y)=[2-C(1-2r)]_+$.
On $3/4\le r\le1$,
$E_C(r,y)\ge C(2r-1)$.
Heat convolution preserves
spatial constants, so
$\mathcal E_C(1,x)\ge
C\int_{3/4}^1(2r-1)\,dr=3C/16$.
\end{proof}

Theorem~\ref{thm:penaltycriterion} is a
zero-compatible \emph{equivalent
criterion}, not a proof of the
required theta-specific inequality.
The genuine missing step is to
establish from the theta correlation
measure, without assuming the sign
of $U_n$, that
\begin{equation}
\boxed{\lim_{C\to\infty}
\mathcal E_{n,C}(T,x)=0
\quad(n\ge1,\ x\in\mathbb R).}
\label{eq:penalty-missing}
\end{equation}
A finite range of Fourier heights
or Laguerre orders is insufficient.
The proved high-order Duhamel
budget saturation also rules out
any order-uniform fractional slack.
The exact-rational model checks
are reproducible in
\path{certificates/xi_penalty_defect_limit_exact.py}
and the algebraic Lean gates in
\path{XiKernel/PenaltyDefectLimit.lean}.
The semigroup, theta-analyticity,
dominated-convergence and
all-order equivalence steps
have not been formalized in Lean;
RH remains unresolved.
